% Cuerpo científico recuperado y distribuido en su residencia terminal.
% Cuerpo editor-ready, sin preámbulo y sin portada autónoma.
% Residencia propuesta: inmediatamente después del teorema del doble círculo
% y antes de Navier--Stokes. No se consume como sustituto de RH ni del
% desarrollo íntegro del doble círculo.

\section{Exact entanglement of nonadic holonomy and spectral phase}
\label{sec:cpe-entrelazamiento-exacto}

Let \(\mathcal H_W\) be the Hilbert space obtained from the positive
Guinand--Weil form, let \(H_W\) be the self-adjoint generator of the
translation group, and let
\[
 \mathcal C_W=(H_W+i/2)(H_W-i/2)^{-1}
\]
its Cayley transform. For each depth \(k\), let \(Z_k\) be the
set of enriched cylindrical states, let \(\mu_k\) be its
projective measure of full support, and let
\(\tau_{k+1,k}:Z_{k+1}\to Z_k\) be the truncation. One considers
\[
 \mathscr K_k=\ell^2(Z_k)\widehat\otimes\mathcal H_W,
 \qquad
 \Omega_k=\sum_{z\in Z_k}\sqrt{\mu_k(z)}\,e_z,
 \qquad
 J_k\psi=\Omega_k\otimes\psi.
\]
If \(z'\) is a descendant of \(z\), one takes
\(w_k(z'\mid z)=\sqrt{\mu_{k+1}(z')/\mu_k(z)}\). For a route
\(\gamma:z\to z'\), the integer memory cocycle is denoted by
\(m(\gamma)\), and the twisted transport is defined by
\begin{equation}
 \widetilde U_k(e_z\otimes\psi)
 :=\sum_{\tau_{k+1,k}z'=z}
 w_k(z'\mid z)\,e_{z'}\otimes
 \mathcal C_W^{\,m(z\to z')}\psi.
 \label{eq:cpe-transporte-cilindrico-cayley}
\end{equation}
The projectivity of \(\mu_k\), the orthogonality of descendants, and the
unitarity of \(\mathcal C_W\) give
\(\widetilde U_k^*\widetilde U_k=I\). The complete label \(z'\)
preserves cylinder, sheet, orientation, carry, route, boundary, and
multiplicity; therefore, distinct histories are not identified upon applying
\eqref{eq:cpe-transporte-cilindrico-cayley}.

\begin{theorem}[Cayley intertwining under \(1\) complete nonadic turn]
\label{thm:cpe-entrelazamiento-cayley-tpk}
If \(\Gamma_9\) is the complete nonadic holonomy, normalized by
\(m(\Gamma_9)=1\), then
\begin{equation}
 \boxed{\widetilde U_{\Gamma_9,k}J_k
 =J_{k+9}\mathcal C_W.}
 \label{eq:cpe-entrelazamiento-exacto}
\end{equation}
For every spectral atom \(\psi_\gamma\) of \(H_W\), with
\(H_W\psi_\gamma=\gamma\psi_\gamma\), one has
\begin{equation}
 \widetilde U_{\Gamma_9^n,k}J_k\psi_\gamma
 =\tau_\gamma^nJ_{k+9n}\psi_\gamma,
 \qquad
 \tau_\gamma=\frac{\gamma+i/2}{\gamma-i/2}\in S^1.
 \label{eq:cpe-orbita-espectral-ledger}
\end{equation}
In the contracting publication, the same identity takes the form
\[
 \widetilde M_{\Gamma_9^n,k}J_k\psi_\gamma
 =\left(\frac59\right)^n\tau_\gamma^n
 J_{k+9n}\psi_\gamma.
\]
\end{theorem}

\begin{proof}
The turn \(\Gamma_9\) returns to the visible phase but increments memory by
one unit, so that \(m(\Gamma_9)=1\). Summing over all the
descendants of \(\Omega_k\), the weights telescope with the projective measure
and yield \(\Omega_{k+9}\), while the Weil fiber receives exactly
one power of \(\mathcal C_W\). This proves
\(\widetilde U_{\Gamma_9,k}J_k=J_{k+9}\mathcal C_W\). The iteration uses the
additivity \(m(\Gamma_9^n)=n\). The final formula follows from
\(\widetilde M_{\Gamma_9,k}=(5/9)\widetilde U_{\Gamma_9,k}\).
\end{proof}

The theorem does not identify a zero with an isolated phase of the calendar. One
complete circuit applies the Cayley character once; each zero determines
the angular law according to which a history traverses all circuits. In
particular, phase return and state return remain distinct:
\[
 g(K+9)=g(K),
 \qquad B_{K+9}\ne B_K\quad\text{in general}.
\]

\section{Convergence in norm of finite spectral matrices}
\label{sec:cpe-convergencia-resolvente}

Let
\[
 R_W=(H_W-iI)^{-1}.
\]
The spectrum of \(H_W\) is discrete and \(|\gamma|\to\infty\), so
\(R_W\) is compact. A countable family
\(\{f_j\}_{j\ge1}\subset C_c^\infty(\mathbb R)\) with rational parameters is chosen,
dense in the test core and selected without consulting spectral heights.
After quotienting by the kernel of the Weil form and orthonormalizing, one
obtains subspaces
\[
 \mathcal V_N=\operatorname{span}\{[f_1],\ldots,[f_N]\},
 \qquad P_N:\mathcal H_W\to\mathcal V_N,
\]
with \(P_N\to I\) strongly. Define
\begin{equation}
 R_N=P_NR_WP_N.
 \label{eq:cpe-resolvente-finito}
\end{equation}

\begin{theorem}[Compact approximation without introducing the zeros]
\label{thm:cpe-convergencia-norma}
The sequence \(R_N\) converges in norm to \(R_W\):
\begin{equation}
 \boxed{\lim_{N\to\infty}\|R_W-R_N\|=0.}
 \label{eq:cpe-convergencia-norma}
\end{equation}
Consequently, every nonzero eigenvalue
\(\lambda_\gamma=(\gamma-i)^{-1}\) of \(R_W\) is the limit, with its
algebraic multiplicity, of eigenvalues of the finite matrices of
\(R_N\). The heights are recovered through
\begin{equation}
 \gamma=i+\lambda_\gamma^{-1}.
 \label{eq:cpe-recuperacion-altura}
\end{equation}
\end{theorem}

\begin{proof}
If \(K\) is compact and \(P_N\to I\) strongly, the convergence is uniform
over the compact \(K(B_1)\), where \(B_1\) is the unit ball. Therefore,
\[
 \|(I-P_N)K\|\longrightarrow0.
\]
Applying the same argument to \(K^*\) yields
\(\|K(I-P_N)\|\to0\). Thus,
\[
 \|K-P_NKP_N\|
 \leq\|(I-P_N)K\|+\|P_NK(I-P_N)\|\longrightarrow0.
\]
Taking \(K=R_W\) proves \eqref{eq:cpe-convergencia-norma}. The
stability of isolated eigenvalues of compact operators under
norm perturbations and the Riesz functional calculus preserve their
multiplicity. The identity \eqref{eq:cpe-recuperacion-altura} is the inverse
of the spectral map of the resolvent.
\end{proof}

The entries of \(R_N\) are computed from the prime--gamma--polar form:
\[
 \langle R_W[f_j],[f_k]\rangle_W
 =i\int_0^\infty e^{-t}
 \mathscr I_{\rm GW}(T_{-t}f_j,f_k)\,dt.
\]
Therefore, the theorem does not use a zero table to form the matrices.

\subsection{A posteriori certification of intervals and multiplicities}

Let \((\lambda_N,v_N)\) be a normalized eigenvector of \(R_N\), and let
\[
 \varepsilon_N(v_N)
 :=\|(I-P_N)R_Wv_N\|.
\]
As \(R_W\) is normal,
\begin{equation}
 \operatorname{dist}(\lambda_N,\sigma(R_W))
 \leq\varepsilon_N(v_N).
 \label{eq:cpe-residuo-inclusion}
\end{equation}
The prime, gamma, and polar channels provide separate interval bounds
for \(\varepsilon_N\); their tails are summed before the interval is emitted.
Moreover, if \(\Gamma\) is a closed contour contained in the resolvent set of
\(R_N\) and
\begin{equation}
 \sup_{z\in\Gamma}
 \|(z-R_N)^{-1}\|\,\|R_W-R_N\|<1,
 \label{eq:cpe-criterio-riesz}
\end{equation}
then the Riesz projections of \(R_N\) and \(R_W\) delimited by
\(\Gamma\) have the same rank. The equations
\eqref{eq:cpe-residuo-inclusion} and \eqref{eq:cpe-criterio-riesz} certify,
respectively, inclusion and completeness of the atoms contained in each
window. Interval arithmetic makes the decimal bounds rigorous; it does not supply a
missing mathematical map.

\section{Coupled trace of the Weil spectrum and the Moonshine grading}
\label{sec:cpe-traza-rh-moonshine}

The product trace
\[
 \operatorname{Tr}_{\mathcal H_W}
 \bigl(e^{-aH_W^2}\mathcal C_W^n\bigr)
 \operatorname{Tr}_{V^\natural}
 \bigl(g e^{-\beta(L_0-1)}\bigr)
\]
superposes two observables, but does not couple them. The memory of the TPK and the
grading of \(V^\natural\) make it possible to construct the coupling without postulating
an action of the Monster on digits or on heights.

Once the exceptional gauge \(\chi_{\rm exc}\) has been fixed, the memory origin and the grading origin are jointly oriented. Let
\(\mathcal H_{\rm mem}=\ell^2(\mathbb N_0)\), with
\(N_{\rm mem}|n\rangle=n|n\rangle\), and let
\[
 V^\natural=\widehat\bigoplus_{r\ge0}V^\natural_r,
 \qquad L_0|_{V^\natural_r}=rI.
\]
Let us denote by \(P_r^\natural\) the projection onto
\(V^\natural_r\). The degree-equality projector is
\begin{equation}
 \Pi_{\chi_{\rm exc}}
 :=\sum_{n\ge0}|n\rangle\langle n|
 \otimes I_{\mathcal H_W}\otimes P_n^\natural.
 \label{eq:cpe-proyector-diagonal-grados}
\end{equation}
The sum converges strongly and defines an orthogonal projection. On the product, the memory character is defined.
\begin{equation}
 \mathcal C_W^{N_{\rm mem}}
 :=\sum_{n\ge0}|n\rangle\langle n|
 \otimes\mathcal C_W^n.
 \label{eq:cpe-cayley-memoria}
\end{equation}

\begin{definition}[Diagonal trace RH--Moonshine]
\label{def:cpe-traza-diagonal}
For \(a,b,\beta>0\) and \(g\in\mathbb M\), one defines
\begin{equation}
\begin{aligned}
 \Theta^{\Delta}_{g,\chi_{\rm exc}}(a,b,\beta)
 :={}&\operatorname{Tr}\!\left[
 \Pi_{\chi_{\rm exc}}
 \Bigl(
 e^{-bN_{\rm mem}}
 \mathcal C_W^{N_{\rm mem}}e^{-aH_W^2}
 \otimes g e^{-\beta(L_0-1)}
 \Bigr)
 \Pi_{\chi_{\rm exc}}
 \right].
 \label{eq:cpe-traza-diagonal-def}
\end{aligned}
\end{equation}
\end{definition}

\begin{theorem}[Existence and Non-Factorized Expansion]
\label{thm:cpe-traza-diagonal}
The trace \eqref{eq:cpe-traza-diagonal-def} is absolutely convergent and
satisfies
\begin{equation}
 \boxed{\Theta^{\Delta}_{g,\chi_{\rm exc}}(a,b,\beta)
 =\sum_{n\ge0}e^{-bn}e^{-\beta(n-1)}
 \operatorname{Tr}_{\mathcal H_W}
 \bigl(e^{-aH_W^2}\mathcal C_W^n\bigr)
 \operatorname{Tr}_{V^\natural_n}(g).}
 \label{eq:cpe-traza-diagonal-expansion}
\end{equation}
After spectral recognition,
\begin{equation}
 \Theta^{\Delta}_{g,\chi_{\rm exc}}(a,b,\beta)
 =\sum_{n\ge0}\sum_\gamma
 m_\gamma e^{-a\gamma^2}
 e^{-bn-\beta(n-1)}\tau_\gamma^n
 \operatorname{Tr}_{V^\natural_n}(g).
 \label{eq:cpe-traza-diagonal-espectral}
\end{equation}
This expression is not the product of a spectral sum and a
McKay--Thompson series: the same index \(n\) simultaneously selects the power
of the memory character and the monstrous degree.
\end{theorem}

\begin{proof}
For \(a>0\), the spectral counting law implies
\(\sum_\gamma m_\gamma e^{-a\gamma^2}<\infty\). Each space
\(V^\natural_n\) is finite-dimensional, and the coefficients of the Moonshine
characters grow subexponentially in \(n\). The factor
\(e^{-(b+\beta)n}\) ensures absolute convergence of the joint series.
The expansion is obtained by inserting the spectral resolutions of
\(N_{\rm mem}\), \(H_W\), and \(L_0\) and using
\eqref{eq:cpe-proyector-diagonal-grados}. The presence of the diagonal projector
eliminates terms with distinct degrees and produces a single sum over
\(n\), rather than the two independent sums of the product trace.
\end{proof}

The coupling depends on the gauge \(\chi_{\rm exc}\), which fixes the origin,
orientation, and alignment of the exceptional reading. A conjugate choice
produces a conjugate trace or a declared reindexing; it does not alter the
Riemann--Weil theorem or the intertwining
\eqref{eq:cpe-entrelazamiento-exacto}.

